\section{搜索 + LLM：企业的杀手级应用}

\begin{importantbox}{Grounding with Search}
几乎每个企业用例都在结合 LLM 和搜索：
\begin{itemize}
    \item LLM 的知识有时效性限制——需要搜索获取最新信息
    \item LLM 会产生幻觉——需要搜索做事实验证
    \item LLM 擅长推理但缺乏引用——搜索提供精确的来源引用
\end{itemize}
这使得 RAG（检索增强生成）成为企业 AI 应用的标配架构。
\end{importantbox}

从工程角度看，RAG pipeline 需要处理三个关键问题：一是向量索引增长带来的 latency，二是检索异常（返回过时或错误文档），三是 LLM 与检索结果的融合策略。Burak 讲到一些团队会在检索结果前加一层「fact filter」——把检索结果先交给 smaller model 进行精排，再交给大型 LLM 生成。

\subsection{本章小结}

搜索 + LLM 是企业最普遍的 GenAI 应用模式，解决了时效性、幻觉和引用三大痛点。

